\subsection*{Supplementary Note S5: benchmark protocol}

This note fixes the rules that every experiment in the Results follows, before any run.
Supplementary Table~\ref{tab:commands} lists the tool versions and the exact commands.
% src: supporting/plan.md section 7 (rules for every run)

\paragraph{Binary.}
Every run uses one tagged impg build, and no run uses the \texttt{target/release} binary that other work rebuilds \authornote{tag name and the owner of the frozen build, decision D6}.
% src: supporting/plan.md sections 4 (binary rebuilt at least 5 times on 2026-11-08) and 9 C1
Next to each result we store \texttt{sha256sum} of the binary, \texttt{impg -{}-version}, the git commit and the tag.
% src: supporting/plan.md section 7
The impg 0.5.0 source under study is commit 0987752, which is 656 commits after the v0.5.0 tag.
% src: git -C /home/guarracino/git/impg describe --tags (v0.5.0-656-g0987752); impg --version (impg 0.5.0)
The build also copies the wfmash and FastGA binaries and the WFA2 shared libraries of wfmash next to impg, so we hash them too.
% src: impg 0987752 build.rs:20-31 (copies wfmash, the WFA2 shared libraries of wfmash, FastGA, FAtoGDB, GIXmake, GIXrm, ALNtoPAF, PAFtoALN, ONEview); VERIFY_CORRECTIONS.md:107

\paragraph{Hardware fields.}
For each node we report the host name, CPU model, sockets, physical cores, hardware threads, RAM in bytes and swap in bytes.
We also report the storage device and file system of inputs and scratch, the kernel, the operating system and the scheduler limits.
The commands \texttt{hostname}, \texttt{lscpu}, \texttt{free -b}, \texttt{df -B1}, \texttt{uname -a} and \texttt{cat /etc/os-release} record these fields.
The workstation has 16 CPUs and 61,478,674,432 bytes of RAM, an AMD Ryzen 7 PRO 7840U, with kernel 6.14.0-37-generic.
% src: supporting/plan.md section 4 (nproc, free -b on 2026-11-08); lscpu and uname -r on 2026-11-09
The cluster node for E2, E3b, E5, E8 and E10 is not chosen \authornote{node and its fields, decision D5 and question Q6}.

\paragraph{Threads.}
The impg option \texttt{-t} sets the thread count, with a default of 4.
% src: impg index --help, impg query --help (-t default 4)
Per-query latency runs use one thread.
Index builds outside the thread sweeps use 16 threads on the workstation and 32 threads on the cluster node.
% src: lscpu on 2026-11-09 (1 socket, 8 cores per socket, 2 threads per core); Table S1 (B) E1, E3a, E4, E6 index -t 16; E2, E3b index -t 32
Thread sweeps use 1, 4, 16 and 48 threads in E2 and 1, 8, 16 and 48 threads in E5, the last as in the Shredtools benchmark.
% src: supporting/plan.md section 7 E2 and E5; Table S1 (B) E2 batch and E5 index loops; Shivakumar2026shredtools paper.txt:346-349
The AGC reader opens at most \texttt{IMPG\_AGC\_DECOMPRESSOR\_SLOTS} decompressors, 8 by default and never more than the thread count.
% src: impg src/agc_index.rs:61-82
A code comment puts each slot at 3.70 GB on the HPRC release 2 archive, so we set the slots to fit free memory.
% src: impg src/agc_index.rs:87-100 (code comment, not a measurement in this paper)
On the workstation we set 2 slots, 7.40 GB by that comment, against 15,510,704,128 available bytes on 2026-11-08.
% src: supporting/plan.md section 7 (free -b on 2026-11-08); 2 x 3.70 = 7.40 GB (src/agc_index.rs:87-100)

\paragraph{Replicates and cache.}
Each per-query measurement in E2 runs 5 times, and we report the median and the 95th percentile over regions and replicates.
% src: supporting/plan.md section 7 E2
Each index build runs 3 times on the workstation \authornote{replicates of the E2 and E5 cluster builds, 1 or 3}.
A warm-cache run follows one untimed run of the same command.
A cold-cache run follows \texttt{sync} and a write of 3 to \texttt{/proc/sys/vm/drop\_caches}, which needs root.
We run cold-cache arms only where the cluster allows this, and we state which arms are warm.
% src: supporting/plan.md section 12 Q6

\paragraph{Resource measurement.}
GNU \texttt{/usr/bin/time -v} wraps every timed command.
We keep the fields Elapsed (wall clock) time, User time, System time, Maximum resident set size and Exit status.
A run counts only with exit status 0, a non-empty output and a row count that the harness checks.
% src: supporting/plan.md section 9 C14
Identity tests compare the sha256 digests of sorted outputs.

\paragraph{Region sampler.}
\texttt{bedtools random} v2.31.1 draws each region set with \texttt{-l} for length, \texttt{-n} for count and a fixed \texttt{-seed} \citep{Quinlan2010bedtools}.
% src: bedtools random -h (v2.31.1, options -l, -n, -seed, -g)
For E1, E2 and E3a it draws 1,200 regions per set, and we split them in the order drawn into 1,000 for tuning and 200 as holdout.
% src: supporting/plan.md section 7 E1; Table S1 (B) E1, E2, E3a blocks (-n 1200); E3b -n 100 per haplotype, E4 and E5 -n 1000, E6 -n 100
We sort each set only after the split.
The genome file comes from \texttt{impg stats -{}-list-sequences}, without its header line, except in the ScRAP arm of E4, which reads the FASTA indexes.
% src: impg src/main.rs:13267-13277 (prints a Sequence and Length header, then name and length); Table S1 (B) E4 block (y.genome from scrap142/*.fa.gz.fai)
The seeds are 1 for E1, 2 for E2, 3 for E3a, 4 for E3b, 5 for E4 and 6 for E5.
They are 7 for the E6 loci, 8 for the E6 C4 haplotypes, 9 for E7 and 10 for E8.
% src: supporting/plan.md section 7 E1 to E8
The harness also fixes the seed of every random pair subset, because \texttt{random:<f>} in impg keeps pairs by a hash threshold.
% src: sweepga 4713ce0 src/knn_graph.rs:371-392

\paragraph{Tuning and held-out sets.}
Tuning sets are the 1,000-region sets, the HPRC year-1 chr19 alignments for E4 and the yeast chrV fixture for E6.
Held-out sets are the 200-region sets, ScRAP 142 for E4 and E9, and the C4 region in 20 haplotypes for E6 \citep{ODonnell2023scrap}.
% src: supporting/plan.md section 7 E1, E4, E6, E9
Only tuning sets may choose \texttt{-d}, \texttt{-m}, \texttt{-{}-min-transitive-len}, \texttt{-{}-min-distance-between-ranges} or subset fractions.
% src: supporting/plan.md section 7 (rules for every run)
We report tuning and held-out scores side by side.
The authors fix the pass targets before the runs \authornote{targets for F1 to F5 and the E8 Jaccard target, decision D2}.

\paragraph{Index mode.}
Every command passes \texttt{-i}, which forces a single index.
% src: impg src/main.rs:11277-11291
The same file list goes to index and query, because the index does not record its file order.
% src: supporting/plan.md section 9 C4 (src/main.rs:11778-11798 vs src/impg.rs:1657-1663)
The tracepoint reader writes a \texttt{.tpa.idx} file next to each TPA file when one is missing, so the TPA folder must be writable.
% src: tpa f21b012 src/reader.rs:176-195 (missing or invalid index is rebuilt, then index.save writes <tpa>.idx); local GRCh38 .tpa.idx of 40,256,530 bytes

\clearpage
\refstepcounter{table}\label{tab:commands}
\noindent\textbf{Supplementary Table~\thetable.} Tool versions to pin and the commands for each experiment E1 to E10, without warm-up runs, build replicates and the commands that make inputs named only in comments.
Every impg flag below appears in \texttt{impg <command> -{}-help} of the binary at commit 0987752, built 2026-11-08 22:10:17 -0700.
% src: help texts in phase2/plan_help_*.txt, byte-identical to the live binary help on 2026-11-09 (cmp)
External flags come from the local \texttt{-{}-help} of each tool.
Shredtools, Mumemto, liftOver, paf2chain and gbz-base are not installed, so their commands are not written.
% src: command -v on the workstation, 2026-11-09
Scripts with the \texttt{.py} suffix belong to the benchmark harness of code fix C14 and do not exist yet.
% src: supporting/plan.md sections 7 and 9 C14

\medskip
\noindent\textbf{(A) Versions.} Versions on the workstation on 2026-11-09, to pin on every machine.

\smallskip
\noindent{\small
\begin{tabular}{@{}lll@{}}
\toprule
Tool & Version & Used in \\
\midrule
impg & tag to cut (C1), source studied at commit 0987752 (impg 0.5.0) & all \\
wfmash & wfmash-rs 07ed9a3, vendoring wfmash 9b2a738 & E4, E5, E6 \\
FastGA & fastga-rs 97a7a56 (pinned in Cargo.lock), bundled with impg & E6, \texttt{-o gfa} arm (default aligner) \\
sweepga & commit 4713ce0 (pinned in Cargo.lock) & E4, E5 \\
seqwish & v0.7.11-35-gaab5bee & E6 \\
vg & v1.76.1 & E2, E8 \\
odgi & v0.9.4-82-g34c461ff & E5 \\
bedtools & v2.31.1 & E1 to E6, E8 \\
samtools, bgzip & 1.21 & E3 to E6 \\
pigz & 2.8 & E3, E7 \\
agc & 3.2.1 & archive checks (C9) \\
aws-cli & 2.36.24 & downloads \\
Python & 3.12.12 & harness \\
GNU time & /usr/bin/time (version string UNKNOWN) & all \\
Shredtools & commit to pin (github.com/vikshiv/shredtools) & E2, E5, E10 \\
Mumemto & v1.4.0 & E5 \\
liftOver & to install and pin & E3a \\
\bottomrule
\end{tabular}}
% src: --version of each tool on 2026-11-09; impg 0987752 Cargo.lock (fastga-rs rev 97a7a56299b5a565964123eb907e93c4480c9f32); citation-keys.md (wfmash-rs 07ed9a3 vendors wfmash 9b2a738); Shivakumar2026shredtools paper.txt:637-640 (Mumemto v1.4.0)

\medskip
\noindent\textbf{(B) Commands.} Shared variables, then one block per experiment.

{\scriptsize
\begin{verbatim}
# bash. IMPG is the tagged build: record sha256sum $IMPG and $IMPG --version
IMPG=/path/to/impg-TAG/target/release/impg
D=/home/guarracino/git/impg/data/hprcv2
P38=$D/GCA_000001405.15_GRCh38_no_alt_analysis_set.PanSN.merged.paf.gz
T38=$D/GCA_000001405.15_GRCh38_no_alt_analysis_set.PanSN.merged.edit-distance.128.tpa
AGC=$D/HPRC_r2_assemblies_0.6.1.agc
X=/home/guarracino/git/impg-paper/experiments            # workstation scratch
S=/scratch/aguarracino/impg-paper                        # cluster scratch
R2=s3://human-pangenomics/pangenomes/freeze/release2
T='/usr/bin/time -v'
export IMPG_AGC_DECOMPRESSOR_SLOTS=2       # workstation: slots x 3.70 GB must fit free RAM

# E1. TPA projection against PAF projection (workstation)
mkdir -p $X/e1 && cd $X/e1
$T $IMPG index -a $P38 -i paf.impg -f -t 16 2> paf.index.time
$T $IMPG index -a $T38 -i tpa.impg -f -t 16 2> tpa.index.time
$IMPG stats -a $P38 -i paf.impg --list-sequences > g38.seqs.tsv
awk -F'\t' 'NR>1 && $1 ~ /^GRCh38#0#chr[0-9XY]+$/ {print $1"\t"$2}' g38.seqs.tsv > g38.genome
bedtools random -l 10000 -n 1200 -seed 1 -g g38.genome > r1200.bed
head -n 1000 r1200.bed | sort -k1,1 -k2,2n > tune.bed
tail -n 200 r1200.bed | sort -k1,1 -k2,2n > hold.bed
for b in tune hold; do for m in 1 2 0; do
  $T $IMPG query -a $P38 -i paf.impg -b $b.bed --no-merge -x -m $m -o bedpe \
    -O paf.$b.m$m -t 16 2> paf.$b.m$m.time
  $T $IMPG query -a $T38 -i tpa.impg -b $b.bed --no-merge -x -m $m --sequence-files $AGC -v 2 \
    -o bedpe -O tpa.$b.m$m -t 16 2> tpa.$b.m$m.time        # -v 2 logs dropped hits
  $T $IMPG query -a $T38 -i tpa.impg -b $b.bed --no-merge -x -m $m --approximate \
    -o bedpe -O apx.$b.m$m -t 16 2> apx.$b.m$m.time
  python3 -I bp_metrics.py --mode identity paf.$b.m$m.bedpe tpa.$b.m$m.bedpe > id.$b.m$m.tsv
  python3 -I bp_metrics.py --mode delta tpa.$b.m$m.bedpe apx.$b.m$m.bedpe > apx.$b.m$m.tsv
done; done
# E1b. A.tpa and B.tpa list each other (picked with data/pairs.py)
# ab.bed: bedtools random -l 10000 -n 100 -seed 1 on the sequences of A
$IMPG index -a A.tpa B.tpa -i ab.bi.impg -f -t 16
$IMPG index -a A.tpa B.tpa -i ab.uni.impg -f --unidirectional -t 16
$IMPG query -a A.tpa B.tpa -i ab.bi.impg -b ab.bed -d 0 -x -m 2 --sequence-files $AGC \
  -o bed -O ab.bi -t 16
$IMPG query -a A.tpa B.tpa -i ab.uni.impg --unidirectional -b ab.bed -d 0 -x -m 2 \
  --sequence-files $AGC -o bed -O ab.uni -t 16

# E2. Index and query the whole release (cluster)
mkdir -p $S/e2/tpas $S/e2/q $S/e2/vg && cd $S/e2
aws s3 cp --no-sign-request --recursive $R2/impg/tpas/ tpas/
ls $PWD/tpas/*.tpa | sort > tpa466.list && sha256sum tpas/*.tpa > tpa466.sha256
$T $IMPG index --alignment-list tpa466.list -i r2.impg -t 32 2> r2.index.time
$IMPG stats --alignment-list tpa466.list -i r2.impg > r2.stats
$IMPG stats --alignment-list tpa466.list -i r2.impg --list-sequences > r2.seqs.tsv
awk -F'\t' 'NR>1 && $1 !~ /^(GRCh38|CHM13)#/ {print $1"\t"$2}' r2.seqs.tsv > nonref.genome
awk -F'\t' 'NR>1 && $1 ~ /^GRCh38#0#chr[0-9XY]+$/ {print $1"\t"$2}' r2.seqs.tsv > g38.genome
for L in 1000 10000 100000 1000000; do
  bedtools random -l $L -n 1200 -seed 2 -g nonref.genome > nr.$L.bed; done
bedtools random -l 10000 -n 1200 -seed 2 -g g38.genome > g38.10000.bed
for m in 1 2; do for rep in 1 2 3 4 5; do
  while read -r c b e n rest; do
    $T $IMPG query --alignment-list tpa466.list -i r2.impg -r "$c:$b-$e" -d 0 -x -m $m \
      --sequence-files $AGC -o bedpe -O q/g38.$m.$rep.$n -t 1 2> q/g38.$m.$rep.$n.time
  done < <(head -n 1000 g38.10000.bed)
done; done
for L in 1000 10000 100000 1000000; do for m in 1 2; do for rep in 1 2 3 4 5; do
  while read -r c b e n rest; do
    $T $IMPG query --alignment-list tpa466.list -i r2.impg -r "$c:$b-$e" -d 0 -x -m $m \
      --sequence-files $AGC -o bedpe -O q/ex.$L.$m.$rep.$n -t 1 2> q/ex.$L.$m.$rep.$n.time
    $T $IMPG query --alignment-list tpa466.list -i r2.impg -r "$c:$b-$e" -d 0 -x -m $m \
      --approximate -o bedpe -O q/ap.$L.$m.$rep.$n -t 1 2> q/ap.$L.$m.$rep.$n.time
  done < <(head -n 1000 nr.$L.bed)
done; done; done
head -n 1000 nr.10000.bed > nr.10000.tune.bed
for t in 1 4 16 48; do
  $T $IMPG query --alignment-list tpa466.list -i r2.impg -b nr.10000.tune.bed -d 0 -x -m 2 \
    --sequence-files $AGC -o bedpe -O batch.t$t -t $t 2> batch.t$t.time; done
$T $IMPG query --alignment-list tpa466.list -i r2.impg -r 'CHM13#0#chr1' -d 0 -o bed \
  --sequence-files $AGC -O chr1 -t 16 2> chr1.time
GBZ=20250930_hprc25272.p98-k311.tmp.fix.v3.gbz
aws s3 cp --no-sign-request $R2/pggb/gfas/whole-genome/$GBZ .
for rep in 1 2 3 4 5; do while read -r c b e n rest; do   # vg chunk -p: 0-based inclusive
  $T vg chunk -x $GBZ -p "$c:$b-$((e-1))" -c 0 -O gfa -t 1 > vg/$n.gfa 2> vg/$n.$rep.time
done < <(head -n 1000 g38.10000.bed); done
aws s3 cp --no-sign-request s3://genome-idx/mumemto/hprc/hprcv2_merged.bumbl .
aws s3 cp --no-sign-request s3://genome-idx/mumemto/hprc/hprcv2_merged.bumbl.bi .
# Shredtools extract on hprcv2_merged.bumbl with g38.10000.bed: not written (not installed)

# E3a. GRCh38 to CHM13 through haplotypes (workstation)
mkdir -p $X/e3a && cd $X/e3a
aws s3 cp --no-sign-request \
  $R2/impg/pafs/all-vs-1/chm13v2.0_maskedY_rCRS.fa.PanSN.merged.paf.gz chm13.paf.gz
pigz -dc $P38 | awk -F'\t' '$1 !~ /^CHM13#/' | bgzip -@ 8 > g38.noCHM13.paf.gz
pigz -dc chm13.paf.gz | awk -F'\t' '$1 !~ /^GRCh38#/' | bgzip -@ 8 > chm13.noGRCh38.paf.gz
$IMPG index -a $P38 chm13.paf.gz -i full.impg -f -t 16
$IMPG index -a g38.noCHM13.paf.gz chm13.noGRCh38.paf.gz -i held.impg -f -t 16
bedtools random -l 1000 -n 1200 -seed 3 -g ../e1/g38.genome > g38.1k.raw.bed
head -n 1000 g38.1k.raw.bed | sort -k1,1 -k2,2n > g38.1k.bed      # last 200: holdout
$IMPG query -a $P38 chm13.paf.gz -i full.impg -b g38.1k.bed -d 0 -o bedpe -O direct -t 16
for m in 2 3; do
  $IMPG query -a g38.noCHM13.paf.gz chm13.noGRCh38.paf.gz -i held.impg -b g38.1k.bed -d 0 \
    -x -m $m -o bedpe -O held.m$m -t 16; done
# chain arm: liftOver of g38.1k.bed with the public chain: not written (not installed)

# E3b. Held-out partners on the release (cluster, after E2)
mkdir -p $S/e3b && cd $S/e3b       # q10.txt: 10 SAMPLE#HAP (seed 4), Qi.paf.gz and Pi per pair
awk -F'\t' 'NR>1 && $1 !~ /^(GRCh38|CHM13)#/ {print $1}' ../e2/r2.seqs.tsv > nonref_contigs.txt
while read -r q; do
  awk -F'\t' -v q="$q#" 'NR>1 && index($1,q)==1 {print $1"\t"$2}' ../e2/r2.seqs.tsv > $q.genome
  bedtools random -l 10000 -n 100 -seed 4 -g $q.genome          # last 20 per Q_i: holdout
done < q10.txt > q10.bed
# per Q_i: drop the partner's records from the per-target PAF of Q_i
pigz -dc Qi.paf.gz | awk -F'\t' -v p="$Pi#" 'index($1,p)!=1' | bgzip -@ 8 > Qi.held.paf.gz
grep -v -F -f q10_tpa_stems.txt ../e2/tpa466.list > held.list
ls $PWD/*.held.paf.gz >> held.list                     # 456 TPA + 10 filtered PAF
$T $IMPG index --alignment-list held.list -i held.impg -f -t 32 2> held.index.time
$IMPG query --alignment-list ../e2/tpa466.list -i ../e2/r2.impg -b q10.bed -d 0 \
  --sequence-files $AGC -o bedpe -O truth.direct -t 16
for m in 1 2 3 4 0; do
  $IMPG query --alignment-list held.list -i held.impg -b q10.bed -d 0 -x -m $m \
    --sequence-files $AGC -o bedpe -O held.refs.m$m -t 16
  $IMPG query --alignment-list held.list -i held.impg -b q10.bed -d 0 -x -m $m \
    --subset-sequence-list nonref_contigs.txt --sequence-files $AGC \
    -o bedpe -O held.norefs.m$m -t 16
done

# E4. Sparse alignment sets (ScRAP 142 holdout here, chr19 tuning after E5)
mkdir -p $X/e4 && cd $X/e4        # scrap142/*.fa.gz: one bgzipped, faidx-indexed FASTA per strain
$IMPG align --sequence-files scrap142/*.fa.gz --wfmash --format joblist --sparsify auto \
  -o jobs_auto                                    # writes jobs_auto/align_jobs.txt
$T $IMPG align --sequence-files scrap142/*.fa.gz --wfmash --format paf --sparsify none \
  -o y_all -t 16 2> y_all.align.time              # writes y_all/alignments.paf
for f in 0.06 0.17; do python3 -I subset_pairs.py --paf y_all/alignments.paf --mode random \
  --fraction $f --seed 1 > r$f.paf; done
python3 -I subset_pairs.py --paf y_all/alignments.paf --mode joblist \
  --jobs jobs_auto/align_jobs.txt > auto.paf
python3 -I subset_pairs.py --paf y_all/alignments.paf --mode star --hub "$HUB" > star.paf
cp y_all/alignments.paf all.paf
for s in all r0.06 r0.17 auto star; do bgzip -f -@ 8 $s.paf
  $IMPG index -a $s.paf.gz -i $s.impg -f -t 16; done
cat scrap142/*.fa.gz.fai | awk -F'\t' -v h="$HUB#" 'index($1,h)!=1 {print $1"\t"$2}' > y.genome
bedtools random -l 10000 -n 1000 -seed 5 -g y.genome | sort -k1,1 -k2,2n > y_r1000.bed
$IMPG query -a all.paf.gz -i all.impg -b y_r1000.bed -d 0 -o bedpe -O truth.direct -t 16
for s in r0.06 r0.17 auto star; do for m in 1 2 3 4 0; do
  $IMPG query -a $s.paf.gz -i $s.impg -b y_r1000.bed -d 0 -x -m $m -o bedpe -O $s.m$m -t 16
done; done

# E5. Gene-level accuracy and Shredtools on HPRC year-1 chr19 (cluster)
mkdir -p $S/e5/hap $S/e5/lat $S/e5/cat && cd $S/e5
aws s3 cp --no-sign-request \
  s3://human-pangenomics/pangenomes/freeze/freeze1/pggb/chroms/chr19.hprc-v1.0-pggb.gfa.gz .
CATB=s3://human-pangenomics/submissions
CATD=FC7E9302-5A35-4A4F-9D4A-B3E957118E36--Y1_CAT_ANNOTATION_GENCODE38
aws s3 cp --no-sign-request --recursive --exclude '*' --include '*.chm13.gff3.gz' \
  $CATB/$CATD/ cat/                                      # 89 files, 8,452,176,760 bytes
zcat chr19.hprc-v1.0-pggb.gfa.gz > chr19.gfa
odgi build -g chr19.gfa -o chr19.og -t 48
odgi paths -i chr19.og -f > chr19.paths.fa
python3 -I split_by_haplotype.py chr19.paths.fa hap/     # bgzip + samtools faidx per haplotype
$IMPG align --sequence-files hap/*.fa.gz --wfmash --format joblist --sparsify auto -o jobs_auto
$T $IMPG align --sequence-files hap/*.fa.gz --wfmash --format paf --sparsify none \
  -o aln_all -t 48 2> align.all.time
cp aln_all/alignments.paf all.paf
python3 -I cat_to_bed.py --gff cat/CHM13.chm13.gff3.gz --chrom chr19 --names seqmap.tsv \
  > cat_genes_chm13_chr19.bed                       # CAT seqid to PanSN path name join
for f in 0.17 0.06; do python3 -I subset_pairs.py --paf all.paf --mode random --fraction $f \
  --seed 1 > r$f.paf; done
python3 -I subset_pairs.py --paf all.paf --mode joblist --jobs jobs_auto/align_jobs.txt > auto.paf
python3 -I subset_pairs.py --paf all.paf --mode star --hub 'CHM13#' > star.paf
for s in all r0.17 r0.06 auto star; do bgzip -f -@ 16 $s.paf
  for t in 1 8 16 48; do
    $T $IMPG index -a $s.paf.gz -i $s.t$t.impg -f -t $t 2> $s.index.t$t.time; done
  [ -s chr19.genome ] || $IMPG stats -a $s.paf.gz -i $s.t48.impg --list-sequences \
    | awk -F'\t' 'NR>1 {print $1"\t"$2}' > chr19.genome
  [ -s rand1000.bed ] || bedtools random -l 10000 -n 1000 -seed 6 -g chr19.genome > rand1000.bed
  $IMPG query -a $s.paf.gz -i $s.t48.impg -b cat_genes_chm13_chr19.bed -d 0 -x -m 2 \
    -o bedpe -O genes.$s -t 16
  while read -r c b e n rest; do
    $T $IMPG query -a $s.paf.gz -i $s.t48.impg -r "$c:$b-$e" -d 0 -x -m 2 -o bed \
      -O lat/$s.$n -t 1 2> lat/$s.$n.time
  done < rand1000.bed
done
python3 -I bp_metrics.py --mode genes --truth cat/ --calls genes.all.bedpe   # each subset
# Shredtools and Mumemto on hap/*.fa.gz: not written (not installed)
# E4 chr19 tuning arm (cluster, E5 folder)
bedtools random -l 10000 -n 1000 -seed 5 -g chr19.genome | sort -k1,1 -k2,2n > c19_r1000.bed
$IMPG query -a all.paf.gz -i all.t48.impg -b c19_r1000.bed -d 0 -o bedpe -O c19.truth -t 16
for s in r0.17 r0.06 auto star; do for m in 1 2 3 4 0; do
  $IMPG query -a $s.paf.gz -i $s.t48.impg -b c19_r1000.bed -d 0 -x -m $m -o bedpe \
    -O c19.$s.m$m -t 16; done; done

# E6. Local graphs and restriction (workstation)
mkdir -p $X/e6 && cd $X/e6
zcat /home/guarracino/git/impg/tests/test_data/yeast.chrV.fa.gz > yeast.chrV.fa  # gzip, not bgzip
samtools faidx yeast.chrV.fa && Y=$PWD/yeast.chrV.fa
$IMPG align --sequence-files $Y --wfmash --format paf --sparsify none -o aln -t 16
seqwish -p aln/alignments.paf -s $Y -g global.k0.gfa -k 0 -t 16
seqwish -p aln/alignments.paf -s $Y -g global.k23.gfa -k 23 -t 16
bgzip -c aln/alignments.paf > aln.paf.gz && $IMPG index -a aln.paf.gz -i aln.impg -f -t 16
while read -r c b e n rest; do                     # loci100.bed: bedtools random -n 100 -seed 7
  $IMPG query -a aln.paf.gz -i aln.impg -r "$c:$b-$e" -x -m 0 --min-transitive-len 0 \
    --min-distance-between-ranges 0 --no-merge -o paf -O C.$n -t 16
  seqwish -p C.$n.paf -s $Y -g R.$n.k0.gfa -k 0 -t 16
  seqwish -p C.$n.paf -s $Y -g R.$n.k23.gfa -k 23 -t 16
  $IMPG query -a aln.paf.gz -i aln.impg -r "$c:$b-$e" -x --no-merge -o paf -O D.$n -t 16
  seqwish -p D.$n.paf -s $Y -g RD.$n.k0.gfa -k 0 -t 16
  $IMPG query -a aln.paf.gz -i aln.impg -r "$c:$b-$e" -x -m 0 --min-transitive-len 0 \
    --min-distance-between-ranges 0 -d 0 -o gfa --gfa-engine seqwish --sequence-files $Y \
    -O T.$n -t 16
  for a in "k0 R.$n.k0" "k23 R.$n.k23" "k0 RD.$n.k0" "k0 T.$n"; do set -- $a
    python3 -I compare_classes.py --global global.$1.gfa --local $2.gfa --closure C.$n.paf \
      > cmp.$2.tsv; done                        # denominator: the bases of closure C
done < loci100.bed
# C4 holdout: 20 haplotypes (seed 8) from the E1 star index; then the same steps on c4_20.fa
$IMPG query -a $P38 -i ../e1/paf.impg -r 'GRCh38#0#chr6:31891045-32123783' -d 0 \
  --subset-sequence-list c4_20_contigs.txt --sequence-files $AGC -o fasta -O c4_20 -t 16
# c4_20.fa names carry ':start-end', so -r on them needs an explicit range, e.g.
# -r 'GRCh38#0#chr6:31891045-32123783:0-232738'

# E7. Copy representation in the released alignments (workstation)
mkdir -p $X/e7 && cd $X/e7          # a1..a3.paf.gz: 3 assembly per-target PAF.gz files, seed 9
for f in $P38 ../e3a/chm13.paf.gz a1.paf.gz a2.paf.gz a3.paf.gz; do
  k=$(basename $f .paf.gz)
  while read -r c b e n rest; do
    pigz -dc -p 4 "$f" | cut -f1-12 \
      | awk -F'\t' -v c="$c" -v b="$b" -v e="$e" '$6==c && $8<e && $9>b' > $k.$n.tsv
    python3 -I locus_depth.py --start $b --end $e $k.$n.tsv >> depth.tsv
  done < loci.$k.bed                                # locus coordinates in each target
done

# E8. Consistency with the released PGGB graph (cluster, after E2)
cd $S/e2 && mkdir -p cons        # cons.bed: 1,000 GRCh38 + 1,000 non-reference 10 kb, seed 10
$IMPG query --alignment-list tpa466.list -i r2.impg -b cons.bed -d 0 -x -m 2 \
  --sequence-files $AGC -o bed -O cons.impg -t 16
while read -r c b e n rest; do
  vg chunk -x $GBZ -p "$c:$b-$((e-1))" -c 0 -O gfa -t 1 > cons/$n.gfa
done < cons.bed
python3 -I chunk_paths_to_bed.py cons/ > cons.vg.bed
python3 -I bp_metrics.py --mode jaccard cons.impg.bed cons.vg.bed --strata strata.bed

# E9. Alignment-backed partition audit on ScRAP 142 (workstation, after E4)
cd $X/e4
$T $IMPG partition -a all.paf.gz -i all.impg -w 10000 -d 1000 --selection-mode haplotype \
  -o bed --output-folder part -t 16 2> part.time
python3 -I exact_once_audit.py --bed-dir part --fai scrap142/*.fa.gz.fai > part.audit.tsv

# E10. Locus refinement against Shredtools zoom (cluster, after E2)
cd $S/e2       # tert_igk.bed: IGKV6-21 CHM13#0#chr2 89168753 89169041, TERT from the annotation
$T $IMPG refine --alignment-list tpa466.list -i r2.impg -b tert_igk.bed -d 10k -x -m 2 \
  --pansn-mode haplotype --span-bp 1000 --max-extension 0.5 --extension-step 1000 \
  --support-output refine.support.bed --sequence-files $AGC -t 16 > refine.tsv 2> refine.time
$T $IMPG refine --alignment-list tpa466.list -i r2.impg -b tert_igk.bed -d 10k -x -m 2 \
  --pansn-mode haplotype --span-bp 1000 --max-extension 0.5 --extension-step 1000 \
  --subset-sequence-list t2t15.txt --support-output refine15.support.bed \
  --sequence-files $AGC -t 16 > refine15.tsv 2> refine15.time
# Shredtools enhance and zoom on hprcv2_merged.bumbl: not written (not installed)
\end{verbatim}
}
